\documentclass[11pt]{amsart}

\usepackage{amsmath,amssymb,amsthm}
\usepackage[margin=1in]{geometry}
\usepackage{hyperref}
\usepackage{enumitem}

\theoremstyle{plain}
\newtheorem{theorem}{Theorem}[section]
\newtheorem{proposition}[theorem]{Proposition}
\newtheorem{lemma}[theorem]{Lemma}
\newtheorem{corollary}[theorem]{Corollary}

\theoremstyle{definition}
\newtheorem{definition}[theorem]{Definition}
\newtheorem{remark}[theorem]{Remark}

\DeclareMathOperator{\Zm}{\mathbb{Z}/m}
\newcommand{\Sm}{S_m(k,\ell)}

\title{A uniform closed form for modular Schur numbers $S_m(k,\ell)$}

\author{Adam McKenna\thanks{\texttt{adam@mysticflounder.ai} \,\textperiodcentered\, \url{https://linkedin.com/in/admckenna} \,\textperiodcentered\, \url{https://mysticflounder.ai}. The Lean~4 proofs of Theorems~\ref{thm:main},~\ref{thm:unified-upper}, and~\ref{thm:unified-lower} were written with Claude (Anthropic) and machine-verified by Aristotle (Harmonic)~\cite{aristotle}.}}
\date{\today}

\begin{document}
\maketitle

\begin{abstract}
The modular Schur number $S_m(k,\ell)$ of Chappelon--Marchena--Dom\'\i nguez~\cite{CMD2013} is the greatest $N$ such that $[1,N]$ admits a $k$-partition into classes that are $\ell$-sum-free modulo~$m$. We prove that
\[
S_m(k,\ell) \;=\; \frac{m}{\gcd(m,\ell-1)} - 1
\]
for every $m\ge 2$, $\ell\ge 2$, and $k\ge m/\gcd(m,\ell-1) - 1$. The proof is two paragraphs of modular arithmetic. It subsumes the per-modulus tables computed in~\cite{CMD2013} ($m\in\{1,2,3\}$) and~\cite{DSWH2025} ($m\in\{4,5,6,7\}$), and extends them to all $m$. The main theorem is formalized in Lean~4 with axiom dependencies only \texttt{propext}, \texttt{Classical.choice}, \texttt{Quot.sound}. We also give a sharp $k_0$ for prime moduli and record a small composite counterexample ($C=\{1,5\}\subseteq\Zm[12]$) that refutes the natural ``only singletons are safe'' conjecture.
\end{abstract}

\section{Introduction}

The classical Schur number $S(k)$ is the largest $N$ such that $\{1,\ldots,N\}$ admits a $k$-partition into sum-free sets. Values are known only through $S(5) = 160\,537$~\cite{Heule2017}; $S(6)$ is the content of Erd\H{o}s problem~483 and remains open. Generalizations to longer equations were studied by Beutelspacher and Brestovansky~\cite{BB1982} and are now textbook material~\cite{LR2014}.

Chappelon, Marchena, and Dom\'\i nguez~\cite{CMD2013} introduced the modular version of the generalized Schur number.

\begin{definition}
A set $S\subseteq\mathbb{Z}$ is \emph{$\ell$-sum-free modulo~$m$} if no $x_1,\ldots,x_\ell\in S$ (repetitions allowed) and $y\in S$ satisfy $x_1+\cdots+x_\ell \equiv y\pmod m$. The \emph{modular Schur number} $S_m(k,\ell)$ is the greatest $N\ge 0$ such that $[1,N]$ admits a $k$-partition into $\ell$-sum-free sets modulo~$m$.
\end{definition}

Chappelon et al.~\cite{CMD2013} closed $m\in\{1,2,3\}$ and observed that the determination of $S_m(k,\ell)$ ``seems to be much more difficult for moduli $m\ge 4$''~(\cite[\S6]{CMD2013}). D'orville, Sim, Wong, and Ho~\cite{DSWH2025} subsequently closed $m\in\{4,5,6,7\}$ via residue-by-residue case analysis on $\ell\bmod m$, leaving $m\ge 8$ as their open problem 1.5. The pattern visible across all ten tables---once $k$ is large enough, $S_m(k,\ell)$ is a function of $\gcd(m,\ell-1)$ alone---was not formulated as a unified statement.

Writing $d:=\gcd(m,\ell-1)$ and $n:=m/d$, the following theorem makes that pattern explicit, with no hypothesis on~$m$ and no case split on $\ell\bmod m$.

\begin{theorem}\label{thm:main}
For every $m\ge 2$, $\ell\ge 2$, and $k\ge n-1$,
\[
S_m(k,\ell) \;=\; n - 1 \;=\; \frac{m}{\gcd(m,\ell-1)} - 1.
\]
\end{theorem}

The proof is elementary: both bounds reduce to the observation that $n$ itself is a singleton that fails $\ell$-sum-freeness. We give it in~\S\ref{sec:upper}--\S\ref{sec:lower}.

The formula looks natural in retrospect, but the prior case-by-case framework obscures it: the proofs in~\cite{CMD2013,DSWH2025} fix a residue class $\ell\bmod m$ and chase a witness partition residue by residue, which hides the single observation that ``the element $n=m/\gcd(m,\ell-1)$ is unsafe as a singleton.'' The argument in \S\ref{sec:upper} is one paragraph because it stops at that observation. The actual combinatorial content is the threshold $k_0$: \S\ref{sec:threshold} gives it tightly for prime $m$ and shows (via the explicit witness $\{1,5\}\subseteq\Zm[12]$) that the prime case does not generalize naively.

\subsection*{Formalization}

The proof is short but has a form that is easy to misstate (an earlier draft confused the role of $n$ versus $m$ in the forbidden tuple). I formalized Theorem~\ref{thm:main} in Lean~4 to remove any doubt; the proof was completed by Aristotle~\cite{aristotle}. The Lean statement
\begin{verbatim}
theorem schurModResidue_eq (m k l : N) (hm : 2 <= m) (hl : 2 <= l)
    (hk : m / Nat.gcd m (l - 1) - 1 <= k) :
    schurModResidue m k l = m / Nat.gcd m (l - 1) - 1
\end{verbatim}
compiles with axiom dependencies only \texttt{propext}, \texttt{Classical.choice}, \texttt{Quot.sound}. Source is at \url{https://github.com/mysticflounder/modular-schur}.

\subsection*{Organization}
\S\ref{sec:prelim} collects base lemmas. \S\ref{sec:upper} and \S\ref{sec:lower} prove Theorem~\ref{thm:main}. \S\ref{sec:threshold} handles the threshold $k_0(m,\ell)$, including the sharp prime-modulus case and a composite counterexample. \S\ref{sec:corollaries} specializes to $m\in\{8,9,11,p^e\}$. \S\ref{sec:open} lists what stays open.

\section{Preliminaries}\label{sec:prelim}

Throughout, $m\ge 2$ and $\ell\ge 2$ are fixed integers, $d := \gcd(m,\ell-1)$, and $n := m/d$. Note $dn = m$ exactly, and $d\mid (\ell-1)$ by definition.

\begin{lemma}[Residue reduction]\label{lem:residues}
If $N\le m-1$, the integers $1,2,\ldots,N$ have pairwise distinct nonzero residues modulo $m$. A partition of $[1,N]$ into classes $C_1\sqcup\cdots\sqcup C_k$ is valid if and only if each $C_i$, viewed as a subset of $\{1,\ldots,m-1\}\subseteq\Zm$, is $\ell$-sum-free modulo $m$.
\end{lemma}
\begin{proof}
All elements of $[1,N]$ lie in $[1,m-1]$, and the map $[1,m-1]\to\Zm$ is injective onto $\{1,\ldots,m-1\}$. Sum-freeness is preserved under this identification.
\end{proof}

\begin{lemma}[Universal upper bound, \cite{CMD2013}]\label{lem:univ-upper}
For every $\ell\ge 1$ and every $k\ge 1$, $\Sm \le m - 1$.
\end{lemma}
\begin{proof}
If $N\ge m$, then $m\in[1,N]$ lies in some class $C$. The tuple $(m,\ldots,m;m)\in C^{\ell+1}$ satisfies $m\cdot\ell\equiv 0\equiv m\pmod m$, violating $\ell$-sum-freeness. (This is~\cite[Eq.~(2)]{CMD2013}.)
\end{proof}

\begin{lemma}[Singleton safety]\label{lem:singleton}
For $r\in\{1,\ldots,m-1\}$, the singleton $\{r\}\subseteq\Zm$ is $\ell$-sum-free modulo $m$ if and only if $(\ell-1)r\not\equiv 0\pmod m$; equivalently, if and only if $n\nmid r$.
\end{lemma}
\begin{proof}
$\ell\{r\}\cap\{r\}\ne\emptyset$ iff $\ell r\equiv r\pmod m$ iff $(\ell-1)r\equiv 0\pmod m$. Write $\ell-1 = du$ so $\gcd(u,n)=1$; then $(\ell-1)r\equiv 0\pmod{dn}$ iff $ur\equiv 0\pmod n$ iff $n\mid r$.
\end{proof}

\begin{lemma}[Monotonicity in $k$]\label{lem:monotone}
If $[1,N]$ admits a valid $k_0$-partition, then for every $k\ge k_0$ it admits a valid $k$-partition.
\end{lemma}
\begin{proof}
Pad with $k-k_0$ empty classes; the empty set is vacuously $\ell$-sum-free.
\end{proof}

\section{Upper bound}\label{sec:upper}

\begin{theorem}\label{thm:unified-upper}
For all $m\ge 2$, $\ell\ge 2$, and $k\ge 1$,
\[
\Sm \;\le\; n - 1.
\]
\end{theorem}
\begin{proof}
Suppose $N\ge n$. Since $1\le n\le m$, we have $n\in[1,N]$; let $C$ be the class of a valid $k$-partition containing $n$. Write $\ell-1 = du$ with $u\in\mathbb{Z}$, so
\[
(\ell - 1)\cdot n \;=\; du\cdot n \;=\; u\cdot(dn) \;=\; u\cdot m \;\equiv\; 0 \pmod m.
\]
Hence $\ell n \equiv n \pmod m$. The tuple $(n,\ldots,n;n)\in C^{\ell+1}$ thus witnesses $\ell$-sum-freeness failure in $C$, a contradiction.
\end{proof}

\begin{remark}
The two obstructions in prior arguments---``$N\ge m$ is blocked by $m$'' (Lemma~\ref{lem:univ-upper}, the $d=1$ case $n=m$) and ``$r=m/d$ is an unsafe singleton'' (Lemma~\ref{lem:singleton} at $r=n$)---are the same obstruction. Both reduce to: the single element $n$ is unsafe as a singleton, so any class containing it is unsafe. Once $N\ge n$, some class contains $n$.
\end{remark}

\section{Lower bound}\label{sec:lower}

\begin{theorem}\label{thm:unified-lower}
For all $m\ge 2$, $\ell\ge 2$, and $k\ge n-1$,
\[
\Sm \;\ge\; n - 1.
\]
\end{theorem}
\begin{proof}
If $n=1$ the conclusion reads $\Sm\ge 0$, vacuously true. Assume $n\ge 2$. Partition $[1,n-1]$ into singletons $D_r := \{r\}$ for $r=1,\ldots,n-1$. By Lemma~\ref{lem:singleton} each $D_r$ is $\ell$-sum-free (since $1\le r<n$ forces $n\nmid r$). This is a valid $(n-1)$-partition; for $k\ge n-1$ pad with empty classes by Lemma~\ref{lem:monotone}.
\end{proof}

\begin{proof}[Proof of Theorem~\ref{thm:main}]
Combine Theorems~\ref{thm:unified-upper} and~\ref{thm:unified-lower}.
\end{proof}

\section{The threshold $k_0(m,\ell)$}\label{sec:threshold}

Let $k_0(m,\ell)$ denote the least $k$ such that $S_m(k',\ell) = n-1$ for every $k'\ge k$.

\begin{theorem}\label{thm:k0-upper}
$k_0(m,\ell) \le \max(1,\,n-1)$.
\end{theorem}
\begin{proof}
Theorem~\ref{thm:unified-lower} exhibits a valid $(n-1)$-partition achieving $N=n-1$ when $n\ge 2$. When $n=1$ the value is zero, achieved trivially at $k=1$.
\end{proof}

\begin{theorem}[Covering lower bound]\label{thm:k0-lower-cov}
Let $\sigma(m,\ell) := \max\{|C| : C\subseteq\{1,\ldots,n-1\},\ C\text{ is $\ell$-sum-free mod }m\}$. Then $k_0(m,\ell)\ge \lceil (n-1)/\sigma(m,\ell)\rceil$.
\end{theorem}
\begin{proof}
If $[1,n-1]$ admits a valid $k$-partition, summing class sizes gives $n-1 \le k\cdot\sigma(m,\ell)$.
\end{proof}

\subsection{Tight threshold for prime modulus}

\begin{theorem}[Large-$\ell$ coset criterion]\label{thm:coset}
Let $C\subseteq\Zm$ with $|C|\ge 2$, pick $a_0\in C$, and let $H := \langle C-C\rangle = g\Zm$ with $g := m/|H|$. For every $\ell\ge |H|-1$,
\[
\ell C \;=\; \ell a_0 + H,
\]
and consequently $C$ is $\ell$-sum-free modulo $m$ if and only if $g\nmid (\ell-1)a_0$.
\end{theorem}
\begin{proof}[Proof sketch]
$(\subseteq)$ Every $\ell$-sum $\sum c_i = \ell a_0 + \sum(c_i-a_0)$ lies in $\ell a_0 + H$ since $c_i-a_0\in H$.

$(\supseteq)$ Pick $a_1\in C$ with $\langle a_1-a_0\rangle = H$ (such $a_1$ exists because differences generate $H$, and in a cyclic group some single difference generates the full subgroup generated by a finite set). Writing $\delta := a_1 - a_0$, the sequence $0,\delta,2\delta,\ldots,(|H|-1)\delta$ exhausts $H$, so for any $h\in H$ there is $j\in\{0,\ldots,|H|-1\}$ with $h = j\delta$; take $c_1=\cdots=c_j = a_1$ and $c_{j+1}=\cdots=c_\ell = a_0$ (valid since $j\le|H|-1\le\ell$) to realize $\ell a_0 + h$ as an $\ell$-sum.

The safety criterion follows: since $C\subseteq a_0+H$, we have $\ell C\cap C\ne\emptyset$ iff $\ell a_0+H = a_0+H$ iff $(\ell-1)a_0\in H$ iff $g\mid(\ell-1)a_0$.
\end{proof}

\begin{remark}
The $(\supseteq)$ direction is a cyclic-group version of Cauchy--Davenport (cf.\ Nathanson~\cite{Nathanson1996}, Ch.~2): over $\Zm[p]$ prime, $\ell$-fold sumsets of any non-degenerate $C$ saturate $\Zm[p]$ once $\ell\ge p-1$. Theorem~\ref{thm:coset} is the saturation statement inside an arbitrary cyclic subgroup $H\le\Zm$, with the explicit $\ell\ge|H|-1$ threshold.
\end{remark}

\begin{corollary}[Sharp $k_0$ for prime moduli]\label{cor:prime-sharp}
Let $p$ be prime. For every $\ell\ge p-1$ with $\ell\not\equiv 1\pmod p$,
\[
k_0(p,\ell) \;=\; p - 1.
\]
\end{corollary}
\begin{proof}
For any $|C|\ge 2$ in $\Zm[p]$, the subgroup $\langle C-C\rangle$ is either $0$ (ruled out) or $\Zm[p]$ (since $\Zm[p]$ has no nontrivial subgroups). By Theorem~\ref{thm:coset} with $g=1$, $\ell C = \Zm[p]$, so $C$ is not $\ell$-sum-free. Hence $\sigma(p,\ell) = 1$, and Theorem~\ref{thm:k0-lower-cov} gives $k_0\ge n-1 = p-1$; Theorem~\ref{thm:k0-upper} supplies the match.
\end{proof}

\subsection{Composite moduli: a counterexample to the natural conjecture}

Corollary~\ref{cor:prime-sharp} might suggest: for every $m$ and every large $\ell$ with $\gcd(m,\ell-1)<m$, only singletons are $\ell$-sum-free, so $k_0(m,\ell) = n-1$. This is false once $m$ is composite.

\begin{proposition}\label{prop:refute-A}
For $m=12$ and every $\ell\equiv 11\pmod{12}$ (so $d=2$, $n=6$), the set $C = \{1,5\}\subseteq\{1,\ldots,5\}$ is $\ell$-sum-free modulo $12$. Consequently $\sigma(12,\ell)\ge 2$ and $k_0(12,\ell)\le 3 < n-1$.
\end{proposition}
\begin{proof}
$C-C = \{0,\pm 4\}$, so $\langle C-C\rangle = 4\mathbb{Z}/12 = \{0,4,8\}$, $|H|=3$, $g=4$. Theorem~\ref{thm:coset} applies once $\ell\ge 2$. With $a_0=1$, the criterion $4\nmid(\ell-1)\cdot 1$ reads $4\nmid\ell-1$, i.e.\ $\ell\not\equiv 1\pmod 4$. For $\ell\equiv 11\pmod{12}$ we have $\ell\equiv 3\pmod 4$, so $4\nmid\ell-1$, and $C$ is $\ell$-sum-free. Combined with singletons $\{2\},\{3\},\{4\}$ this gives a valid $3$-partition of $[1,5]$.
\end{proof}

\section{Corollaries: recovered and extended formulas}\label{sec:corollaries}

Specializing Theorem~\ref{thm:main}:

\begin{itemize}[leftmargin=2em]
  \item \textbf{$m = 8$:} for $k\ge 7$, $S_8(k,\ell)$ equals $0,1,3,7$ according to $\gcd(8,\ell-1)\in\{8,4,2,1\}$ (residue pattern $\ell\bmod 8$).
  \item \textbf{$m = 9$:} for $k\ge 8$, $S_9(k,\ell)$ equals $0,2,8$ according to $\gcd(9,\ell-1)\in\{9,3,1\}$.
  \item \textbf{$m = 11$:} for $k\ge 10$, $S_{11}(k,\ell) = 0$ if $\ell\equiv 1\pmod{11}$, else $10$.
  \item \textbf{Prime-power $m=p^e$:} writing $v:=\min(v_p(\ell-1),e)$, for $k\ge p^{e-v}-1$,
  \[
  S_{p^e}(k,\ell) \;=\; p^{e-v} - 1.
  \]
\end{itemize}
All of these stabilized rows are consistent with (and for $m=4,5,6,7$ recover) the tables of \cite[\S 2--3]{DSWH2025}.

\section{What stays open}\label{sec:open}

\begin{enumerate}[leftmargin=2em,label=(\arabic*)]
  \item \emph{Tight $k_0$ for composite $m$.} The covering bound (Theorem~\ref{thm:k0-lower-cov}) matches the easy upper bound $n-1$ for primes but not in general. For $m=8$, $\ell$ even: $\sigma(8,\ell)=4$ (witness $\{1,3,5,7\}$), so the covering bound gives only $k_0\ge 2$, while the true value is $k_0=3$. A closed form for $k_0(m,\ell)$ at composite $m$ is open.
  \item \emph{Boundary $k < n-1$.} Theorem~\ref{thm:main} only governs $k\ge n-1$. The smaller-$k$ regime---where classes must be large and $\ell$-sum-free---is empirically eventually periodic in $\ell\bmod m$ for fixed $k$ (verified through $m\le 13$), but I have no formula.
  \item \emph{Optimizing over $\ell$.} For fixed $k$ and $m\to\infty$, what is $\max_\ell S_m(k,\ell)$? The pointwise value is Theorem~\ref{thm:main}; the sup over $\ell$ becomes a divisor-counting problem on $m$.
\end{enumerate}

\subsection*{Acknowledgments}

This paper was written collaboratively with Claude (Anthropic), and the Lean~4 proofs were machine-verified by Aristotle (Harmonic)~\cite{aristotle}. The informal proofs, formalization, and draft were iterated together over the course of several sessions. All mathematical errors are mine.

\begin{thebibliography}{9}

\bibitem{BB1982}
A.~Beutelspacher and W.~Brestovansky,
\emph{Generalized Schur numbers},
Lecture Notes in Mathematics \textbf{969} (1982), 30--38.

\bibitem{CMD2013}
J.~Chappelon, M.~P.~Revuelta Marchena, and M.~I.~Sanz Dom\'\i nguez,
\emph{Modular Schur numbers},
Electron. J. Combin. \textbf{20}(2) (2013), \#P61. \texttt{arXiv:1306.5635}.

\bibitem{DSWH2025}
R.~D'orville, C.~Sim, E.~Wong, and T.-K.~L.~Ho,
\emph{Modular generalizations of Schur numbers},
Integers \textbf{25} (2025), \#A62.
\url{https://math.colgate.edu/~integers/z62/z62.pdf}

\bibitem{Heule2017}
M.~J.~H.~Heule,
\emph{Schur number five},
AAAI 2018, 6598--6606. \texttt{arXiv:1711.08076}.

\bibitem{LR2014}
B.~M.~Landman and A.~Robertson,
\emph{Ramsey Theory on the Integers}, 2nd ed.,
Student Mathematical Library \textbf{73}, AMS, 2014.

\bibitem{Nathanson1996}
M.~B.~Nathanson,
\emph{Additive Number Theory: Inverse Problems and the Geometry of Sumsets},
Graduate Texts in Mathematics \textbf{165}, Springer, 1996.

\bibitem{aristotle}
T.~Achim \emph{et al.},
\emph{Aristotle: IMO-level Automated Theorem Proving},
2025. \texttt{arXiv:2510.01346}.

\end{thebibliography}

\end{document}
